\section{The Parallel Approach}\label{sec:approach}

\subsection{The purity of successor computation}

Re-examine \cref{alg:seq}.  The call $\succfn(\ell, Z)$ at line~3 computes a
set of successor symbolic states.  Its inputs are:
\begin{itemize}
  \item The automaton $\mathcal{A}$: a globally shared, \emph{read-only}
        structure after initialisation.
  \item The location $\ell$ and zone $Z$: private to the current iteration
        and never modified by any other iteration.
\end{itemize}
Its output is a freshly allocated vector of \texttt{(loc, DBM)} pairs.
It neither reads nor writes $\Passed$, $\Wait$, or any structure belonging to
any other iteration.

\begin{center}
\fbox{%
\parbox{0.82\linewidth}{%
\textbf{Key observation.}\ $\succfn(\ell, Z)$ is a \emph{pure function}:
given the (read-only) automaton, its return value depends only on $(\ell, Z)$
and is completely independent of the passed list, the waiting queue, and
all other symbolic states.
}}
\end{center}

\noindent
This purity is what makes the calls to $\succfn$ across different states
in the same BFS level \emph{embarrassingly parallel}: they share no
input data that any of them modifies.

\subsection{Why UPPAAL cannot exploit this directly}

In the sequential loop of \cref{alg:seq}, purity of $\succfn$ alone is not
enough: after computing the successors, each iteration immediately checks
them against $\Passed$ and writes new ones back.  The check on line~4 depends
on line~6 of an earlier iteration: a zone written by iteration $i$ might
subsume a zone generated by iteration $i+1$, so the result of $i+1$'s
subsumption test depends on $i$'s write.  Because iterations are pulled off a
FIFO queue, there is no bound on how far apart two dependent iterations can
be---they might belong to entirely different BFS levels, or they might be at
the same BFS level via two different paths to the same location.

In other words: the \emph{per-iteration} structure interleaves pure work
(compute $\succfn$) with stateful work (check and update $\Passed$), and this
interleaving is what forces sequential execution.

\subsection{Level-synchronous parallel BFS: the idea}

The fix is to \emph{separate} the two phases globally rather than
per-iteration.  Instead of processing one state at a time, we collect all
states at the same BFS level into a \emph{frontier array} $\Frontier_k$ and
proceed in two phases per level:

\medskip
\noindent\textbf{Phase~1 (parallel expand).}
Every state $(\ell_i, Z_i) \in \Frontier_k$ is expanded by
\emph{calling $\succfn(\ell_i, Z_i)$ in parallel}.  Thread $t$ appends its
results to a \emph{private} buffer $\mathit{buf}_t$.  No thread reads or
writes $\Passed$.  No synchronisation primitive is acquired.

\medskip
\noindent\textbf{Phase~2 (sequential merge).}
A single thread iterates over all buffers in a fixed order, performing the
same subsumption check-and-insert that \cref{alg:seq} does.  States that
survive subsumption form $\Frontier_{k+1}$.

\medskip
\noindent
\Cref{fig:two-phases} illustrates the contrast between the sequential loop
and the two-phase parallel loop.  \Cref{alg:par} gives the full parallel
algorithm.

\subsection{The parallel algorithm}

\begin{algorithm}[t]
\caption{Level-synchronous parallel BFS (OpenMP backend)}\label{alg:par}
\begin{algorithmic}[1]
\Require TA $\mathcal{A}$,\ target $(\ell^\star, \phi)$,\ thread count $T$
\State $Z_0 \leftarrow \mathit{init}(\mathcal{A})$;\enspace
       $\Passed[\Init] \leftarrow \{Z_0\}$;\enspace
       $\Frontier \leftarrow [(\Init, Z_0)]$
\While{$\Frontier \neq \emptyset$}
  \vspace{2pt}
  \State \textit{// --- Phase 1: embarrassingly parallel expand ---}
  \State $\mathit{buf}_t \leftarrow []$ \textbf{ for } $t = 0,\ldots,T{-}1$
  \ParFor{$i = 0,\ldots,|\Frontier|{-}1$}
    \Comment{\texttt{\#pragma omp parallel for schedule(dynamic,8)}}
    \State $\mathit{buf}_{\mathit{tid}} \mathrel{+}=
           \succfn(\Frontier[i].\ell,\; \Frontier[i].Z)$
    \Comment{writes only to private $\mathit{buf}_{\mathit{tid}}$}
  \EndParFor
  \vspace{2pt}
  \State \textit{// --- Phase 2: sequential merge ---}
  \State $\mathit{next} \leftarrow []$
  \ForAll{$t = 0,\ldots,T{-}1$;\enspace $(\ell', Z') \in \mathit{buf}_t$}
    \If{$\mathit{isTarget}(\ell', Z')$} \Return \textbf{true} \EndIf
    \If{$\exists Z_p \in \Passed[\ell'] : Z' \subseteq Z_p$}
      \textbf{continue}
    \EndIf
    \State $\Passed[\ell'] \mathrel{+}= Z'$;\enspace
           $\mathit{next} \mathrel{+}= (\ell', Z')$
  \EndFor
  \State $\Frontier \leftarrow \mathit{next}$
\EndWhile
\Return \textbf{false}
\end{algorithmic}
\end{algorithm}

\subsection{What changes and what stays the same}

\begin{description}
  \item[Unchanged.] The successor function $\succfn$ is identical to the
    sequential version: same guard intersection, reset, invariant, delay, and
    LU$^+$ extrapolation.  The subsumption predicate ($Z' \subseteq Z_p$) is
    identical.  The passed list structure is identical.

  \item[Changed: data is grouped by level.]  In \cref{alg:seq}, states are
    mixed across all BFS depths in a single FIFO.  In \cref{alg:par}, states
    are explicitly grouped by depth into $\Frontier$.  A state at depth $k$
    can only subsume states at depth $\leq k$, so all states in $\Frontier_k$
    were already checked against $\Passed$ before being placed there---their
    expand calls are safe to run in any order.

  \item[Changed: merge order is thread-then-append.]  The sequential loop
    processes successors immediately as it generates them.  \Cref{alg:par}
    defers all subsumption checks to Phase~2, iterating buffers in thread-ID
    order.  The verdict is the same (see \cref{sec:correctness}), but the
    exact set of zones stored in $\Passed$ may differ in cases where two
    candidates at the same level mutually subsume each other (a rare case
    that does not affect reachability).

  \item[New: thread-local buffers.]  Each thread writes to its own
    \texttt{std::vector<(int,\,Dbm)>}, pre-allocated once per level.  No
    atomic operations, no mutexes, no cache-line sharing between threads
    during Phase~1.
\end{description}

\subsection{Why the sequential merge is not the bottleneck}

A natural concern is whether Phase~2 negates the gains of Phase~1.  In
practice it does not, for two reasons.

\textbf{Bounded passed list.}  LU$^+$ extrapolation ensures that
$|\Passed[\ell]|$ is bounded by the finite LU$^+$-quotient of the zone graph,
a constant that depends only on the clock constants of the model.  For all
benchmarks we have run, this bound is at most a few hundred per location.
The subsumption check for one candidate is therefore $O(n^2)$ amortised.

\textbf{Cheap per-state merge.}  The merge does one hash lookup
(exact-duplicate check in $O(n^2)$ via a pre-computed byte-string key) and at
most $|\Passed[\ell]|$ DBM inclusion tests.  This is orders of magnitude
cheaper than generating the successors (which requires up to $|\delta(\ell)|$
full DBM operations at $O(n^3)$ each).

Consequently, Phase~2 accounts for a small fraction of total wall-clock time,
and Phase~1 scales with $T$.
